\apendice{Trecho da implementação JuMP da formulação F2}
\label{ap:A}

A título de ilustração, reproduz-se a seguir o núcleo do \textit{script} \verb|(F2)QLIM+VLIM.jl|, responsável pela construção do modelo JuMP descrito na \autoref{subsec:f16}. O código completo, junto com as demais formulações da progressão F0 a F5, está disponível no repositório do trabalho em \verb|powerflow/Codes/PF_Formulation/|.

\begin{verbatim}
# Variáveis de estado físico
@variable(model,
    ref[:bus][i]["vmin"] <= vm[i in keys(ref[:bus])] <= ref[:bus][i]["vmax"],
    start = ref[:bus][i]["vm"])
@variable(model, va[i in keys(ref[:bus])], start = ref[:bus][i]["va"])

@variable(model,
    ref[:gen][i]["pmin"] <= pg[i in keys(ref[:gen])] <= ref[:gen][i]["pmax"],
    start = ref[:gen][i]["pg"])
@variable(model,
    ref[:gen][i]["qmin"] <= qg[i in keys(ref[:gen])] <= ref[:gen][i]["qmax"],
    start = ref[:gen][i]["qg"])

# Slacks de controle (VLIM e QLIM)
PENALIDADE = 1e6
gen_buses = unique([gen["gen_bus"] for (i, gen) in ref[:gen]])

@variable(model, sl_v[i in gen_buses], start = 0.0)
for bus_id in gen_buses
    vm_setpoint = ref[:bus][bus_id]["vm"]
    @constraint(model, vm[bus_id] == vm_setpoint + sl_v[bus_id])
end

@variable(model, sl_d[i in keys(ref[:load])], start = 0.0)

# Equações de fluxo nos ramos (modelo pi com transformador)
for (l, branch) in ref[:branch]
    f = branch["f_bus"]; t = branch["t_bus"]
    g, b   = PowerModels.calc_branch_y(branch)
    tr, ti = PowerModels.calc_branch_t(branch)
    tm     = branch["tap"]
    g_fr, b_fr = branch["g_fr"], branch["b_fr"]
    g_to, b_to = branch["g_to"], branch["b_to"]

    p[(l, f, t)] = @NLexpression(model,
        (g + g_fr) / tm^2 * vm[f]^2
        + (-g*tr + b*ti) / tm^2 * (vm[f]*vm[t]*cos(va[f]-va[t]))
        + (-b*tr - g*ti) / tm^2 * (vm[f]*vm[t]*sin(va[f]-va[t])))
    # ... análogo para q[(l,f,t)], p[(l,t,f)], q[(l,t,f)]
end

# Função objetivo soft-constraint
@objective(model, Min,
    PENALIDADE * sum(sl_v[i]^2 for i in keys(sl_v)) +
    PENALIDADE * sum(sl_d[l]^2 for l in keys(sl_d)))

optimize!(model)
\end{verbatim}

A leitura do trecho acima permite identificar três blocos sintáticos centrais à abordagem deste trabalho: (i) declaração de variáveis primárias com limites caixa e valores iniciais provenientes do PWF; (ii) introdução das folgas \verb|sl_v| e \verb|sl_d| associadas a VLIM e QLIM; (iii) escrita das equações de fluxo do ramo na forma do modelo $\pi$ com transformador, usando \verb|@NLexpression| do JuMP para registrar o caráter não linear das expressões. As formulações F3 a F5 acrescentam, sucessivamente, blocos análogos para $b^{sh}_s$ (CSCA), $t_{ij}$ (CTAP) e injeções DERA, sem modificar a estrutura geral.
